% ============================================================================
% BIÊN BẢN HỌP NHÓM SỐ 01 – DỰ ÁN SMART E-MOBILITY HUB
% Môn: CO3001 – Công nghệ Phần mềm (HK261)
% Trường Đại học Bách Khoa TP.HCM – Khoa KH & KT Máy tính
% Compiled with: pdflatex (2 passes recommended)
% ============================================================================

\documentclass[11pt,a4paper]{article}

% --- Ngôn ngữ & Encoding ---
\usepackage[utf8]{vietnam}

% --- Geometry ---
\usepackage[margin=2cm, headheight=22pt, top=2.5cm]{geometry}

% --- Bảng biểu & Hình ảnh ---
\usepackage{booktabs}
\usepackage{tabularx}
\usepackage{array}
\usepackage{multirow}
\usepackage{graphicx}
\usepackage{longtable}
\usepackage{colortbl}

% --- Hyperlinks ---
\usepackage[colorlinks=true, linkcolor=hcmutblue, urlcolor=hcmutblue, citecolor=hcmutblue]{hyperref}
\PassOptionsToPackage{hyphens}{url}
\urlstyle{tt}

% --- Màu sắc ---
\usepackage[dvipsnames,svgnames,x11names]{xcolor}
\definecolor{hcmutblue}{RGB}{0, 83, 159}      % HCMUT Brand Blue
\definecolor{hcmutred}{RGB}{180, 30, 40}       % Accent Red
\definecolor{lightgray}{RGB}{240, 240, 240}
\definecolor{darkgray}{RGB}{80, 80, 80}
\definecolor{sectionblue}{RGB}{20, 60, 120}

% --- Danh sách ---
\usepackage{enumitem}

% --- Header / Footer ---
\usepackage{fancyhdr}
\pagestyle{fancy}
\fancyhf{}
\fancyhead[L]{\footnotesize\textcolor{hcmutblue}{\textbf{HCMUT} -- Khoa Khoa học \& Kỹ thuật Máy tính}}
\fancyhead[R]{\footnotesize\textcolor{hcmutblue}{CO3001 -- Công nghệ Phần mềm (HK261)}}
\fancyfoot[C]{\footnotesize Trang \thepage}
\renewcommand{\headrulewidth}{0.8pt}
\renewcommand{\footrulewidth}{0.4pt}
\renewcommand{\headrule}{\hbox to\headwidth{\color{hcmutblue}\leaders\hrule height \headrulewidth\hfill}}
\renewcommand{\footrule}{\hbox to\headwidth{\color{hcmutblue}\leaders\hrule height \footrulewidth\hfill}}

% --- Định dạng tiêu đề section ---
\usepackage{titlesec}
\titleformat{\section}
  {\Large\bfseries\color{sectionblue}}
  {\thesection.}{0.5em}{}
  [\vspace{-0.5em}\textcolor{hcmutblue}{\rule{\textwidth}{0.6pt}}]

\titleformat{\subsection}
  {\large\bfseries\color{darkgray}}
  {\thesubsection.}{0.5em}{}

% --- Khoảng cách dòng ---
\usepackage{setspace}
\setstretch{1.25}

% --- Loại bỏ indent đoạn, thêm khoảng cách giữa các đoạn ---
\setlength{\parindent}{0pt}
\setlength{\parskip}{6pt}

% --- Tiện ích vẽ đường kẻ ngang ---
\newcommand{\sectionrule}{\vspace{4pt}\noindent\textcolor{hcmutblue}{\rule{\textwidth}{0.4pt}}\vspace{4pt}}

% ============================================================================
\begin{document}

% ============================================================================
%                           KHỐI TIÊU ĐỀ (HEADER BLOCK)
% ============================================================================
\thispagestyle{fancy}

\begin{center}
  {\small\textcolor{darkgray}{TRƯỜNG ĐẠI HỌC BÁCH KHOA TP.HCM -- ĐẠI HỌC QUỐC GIA TP.HCM}}\\[2pt]
  {\small\textcolor{darkgray}{KHOA KHOA HỌC \& KỸ THUẬT MÁY TÍNH}}\\[2pt]
  \textcolor{hcmutblue}{\rule{0.6\textwidth}{1.2pt}}\\[12pt]
  {\LARGE\bfseries\textcolor{hcmutblue}{BIÊN BẢN HỌP NHÓM SỐ 01}}\\[6pt]
  {\Large\textcolor{sectionblue}{DỰ ÁN SMART E-MOBILITY HUB}}\\[4pt]
  {\normalsize\textit{Hệ thống Điều phối Phương tiện Điện -- Khu Đô thị ĐHQG-HCM}}\\[4pt]
  \textcolor{hcmutblue}{\rule{0.6\textwidth}{1.2pt}}
\end{center}

\vspace{8pt}

% ============================================================================
%                          BẢNG THÔNG TIN TỔNG QUÁT
% ============================================================================
\begin{center}
\renewcommand{\arraystretch}{1.4}
\begin{tabularx}{\textwidth}{|>{\bfseries\raggedright\arraybackslash}p{3.8cm}|X|}
  \hline
  \rowcolor{lightgray} \multicolumn{2}{|c|}{\textbf{THÔNG TIN CUỘC HỌP}} \\
  \hline
  Môn học        & CO3001 -- Công nghệ Phần mềm, Học kỳ 1/2026--2027 (HK261) \\
  \hline
  Thời gian      & 20:00 -- 20:30, Thứ Bảy ngày 12/09/2026 \\
  \hline
  Hình thức      & Trực tuyến qua \textbf{Google Meet} \\
  \hline
  Chủ trì        & Nguyễn Thành Đạt (MSSV: \textit{2410709}) \\
  \hline
  Thư ký         & Nguyễn Thành Đạt (MSSV: \textit{2410709}) \\
  \hline
  Thành phần      & Đầy đủ \textbf{7/7} thành viên (không có vắng mặt) \\
  \hline
\end{tabularx}
\end{center}

\vspace{10pt}

% ============================================================================
%               MỤC 1: MỤC TIÊU CUỘC HỌP
% ============================================================================
\section{Mục tiêu cuộc họp}

\begin{enumerate}[label=\textbf{\arabic*.}, leftmargin=2em, itemsep=3pt]
  \item \textbf{Khởi động dự án (Kickoff):} Khai mạc dự án BTL Công nghệ Phần mềm -- Hệ thống Điều phối Phương tiện Điện tại Khu đô thị ĐHQG-HCM (Smart E-Mobility Hub).
  \item \textbf{Phân chia vai trò \& trách nhiệm:} Phân công chính cho từng thành viên trong nhóm phụ trách 7 phân hệ nghiệp vụ (Subsystems) và các Use-case cá nhân.
  \item \textbf{Thống nhất quy chế vận hành:} Cam kết về tiến độ làm việc nhóm, kênh liên lạc, lịch họp định kỳ, deadline nội bộ và chính sách liêm chính học thuật (GenAI Policy).
  \item \textbf{Chốt công nghệ \& công cụ:} Thống nhất Tech Stack (Git/GitHub, Python, Streamlit, Plotly, Draw.io) và kiểm tra khởi tạo repository mã nguồn.
\end{enumerate}

% ============================================================================
%               MỤC 2: QUY CHẾ LÀM VIỆC NHÓM
% ============================================================================
\section{Quy chế làm việc nhóm (Team Operating Regulations)}

\begin{itemize}[leftmargin=2em, itemsep=4pt]
  \item \textbf{Kênh liên lạc chính:} Nhóm chat dự án (Zalo) và GitHub Issues / Pull Requests.
  \item \textbf{Cam kết thời gian phản hồi:} Tất cả thành viên bắt buộc phải phản hồi thông báo/tin nhắn dự án trong vòng tối đa \textbf{04 giờ} trong giờ hành chính (08:00~--~21:00).
  \item \textbf{Lịch họp định kỳ:} Họp tuần trực tuyến vào tối Chủ nhật hàng tuần lúc 20:00 qua Google Meet. Vắng mặt phải báo trước ít nhất 12 giờ kèm lý do chính đáng.
  \item \textbf{Hạn nộp nội bộ (Internal Deadline):} Mọi công việc cá nhân phải hoàn thành sớm hơn hạn nộp chính thức của LMS từ \textbf{24--48 giờ} để thực hiện Peer Review và tổng hợp file.
\end{itemize}

% ============================================================================
%               MỤC 3: PHÂN CÔNG 7 PHÂN HỆ NGHIỆP VỤ
% ============================================================================
\section{Phân công 7 Phân hệ Nghiệp vụ \& Use-case cá nhân (Submission \#1)}

Nhóm đã thống nhất phân công 7 phân hệ chức năng và các Use-case chi tiết tương ứng cho Submission \#1 như bảng dưới đây. Mỗi thành viên chịu trách nhiệm đặc tả chi tiết 02 Use-case cá nhân và thực hiện review chéo (Peer Review) cho thành viên được chỉ định.

\vspace{6pt}

\renewcommand{\arraystretch}{1.45}
{\footnotesize
\begin{tabularx}{\textwidth}{|c|p{3.0cm}|p{3.6cm}|X|p{2.5cm}|}
  \hline
  \rowcolor{lightgray} \textbf{STT} & \textbf{Thành viên \& Vai trò} & \textbf{Phân hệ Nghiệp vụ} & \textbf{Use-case cá nhân đặc tả} & \textbf{Người Review chéo} \\
  \hline
  \textbf{01} &
  \textbf{Nguyễn Thành Đạt}\newline\textit{Leader / PM} &
  \textbf{SS1:} Quản lý \& Đặt xe cho Sinh viên &
  \texttt{UC\_Dat\_Xe\_Dung\_Chung}\newline
  \texttt{UC\_Huy\_Dat\_Xe} &
  Nguyễn Anh Tài \\
  \hline
  \textbf{02} &
  \textbf{Nguyễn Anh Tài}\newline\textit{Dev / Analyst} &
  \textbf{SS2:} Đăng ký dịch vụ Xe cá nhân &
  \texttt{UC\_Dang\_Ky\_Gui\_Xe\_Ca\_Nhan}\newline
  \texttt{UC\_Dat\_Lich\_Sac\_Xe\_Ca\_Nhan} &
  Nguyễn Thành Đạt \\
  \hline
  \textbf{03} &
  \textbf{Phạm Tấn Đạt}\newline\textit{Dev / Analyst} &
  \textbf{SS3:} Giám sát mạng lưới Hub &
  \texttt{UC\_Giam\_Sat\_Trang\_Thai\_Hub}\newline
  \texttt{UC\_Canh\_Bao\_Qua\_Tai} &
  Đặng H. Quý Nhân \\
  \hline
  \textbf{04} &
  \textbf{Đặng H. Quý Nhân}\newline\textit{Dev / Analyst} &
  \textbf{SS4:} Điều phối PTV \& Sự cố &
  \texttt{UC\_Lap\_Lenh\_Dieu\_Chuyen}\newline
  \texttt{UC\_Bao\_Cao\_Su\_Co\_PTV} &
  Phạm Tấn Đạt \\
  \hline
  \textbf{05} &
  \textbf{Dương Đăng Khoa}\newline\textit{Dev / Analyst} &
  \textbf{SS5:} Lập lịch sạc thông minh &
  \texttt{UC\_Lap\_Lich\_Uu\_Tien\_Sac}\newline
  \texttt{UC\_Giam\_Sat\_Hang\_Cho\_Sac} &
  Nguyên \\
  \hline
  \textbf{06} &
  \textbf{Hoàng Văn Tấn}\newline\textit{Dev / Analyst} &
  \textbf{SS6:} Mô phỏng What-if -- Nhu cầu &
  \texttt{UC\_Mo\_Phong\_Tang\_Dot\_Bien}\newline
  \texttt{UC\_Phan\_Tich\_Anh\_Huong} &
  Dương Đăng Khoa \\
  \hline
  \textbf{07} &
  \textbf{Nguyên}\newline\textit{Dev / Analyst} &
  \textbf{SS7:} Mô phỏng What-if -- Sự cố &
  \texttt{UC\_Mo\_Phong\_Su\_Co\_Mat\_Dien}\newline
  \texttt{UC\_Goi\_Y\_Dieu\_Huong} &
  Hoàng Văn Tấn \\
  \hline
\end{tabularx}
}

\vspace{6pt}
{\small\textit{Chi tiết đầy đủ xem tại:} \texttt{docs/tasks-assignment.md}}

% ============================================================================
%               MỤC 4: MỤC TIÊU & YÊU CẦU CẦN ĐẠT SAU SPRINT 1
% ============================================================================
\section{Mục tiêu \& Yêu cầu cần đạt được sau Sprint 1 (13/09 -- 20/09)}

Sprint 1 kéo dài \textbf{01 tuần} (13/09 -- 20/09/2026) tập trung hoàn thành toàn bộ nền tảng phân tích yêu cầu phục vụ Submission \#1. Ba kết quả bắt buộc cần đạt:

\subsection{Understanding Project's Context -- Hiểu rõ bối cảnh dự án}

\begin{itemize}[leftmargin=2em, itemsep=3pt]
  \item Thấu hiểu sâu sắc bài toán điều phối phương tiện điện (Smart E-Mobility Hub) tại Khu đô thị ĐHQG-HCM.
  \item Xác định rõ bối cảnh kết nối \textbf{Ga Metro số 1}, \textbf{KTX Khu A/B}, phạm vi quản lý trạm sạc và bãi đỗ phương tiện.
  \item Phân tích toàn diện kỳ vọng của các bên liên quan (\textit{Stakeholders}):
    \begin{itemize}[leftmargin=1.5em, itemsep=1pt]
      \item \textbf{Sinh viên:} Đặt xe dùng chung, gửi xe cá nhân, tra cứu lịch sạc.
      \item \textbf{Operator:} Giám sát trạng thái hub, điều phối sự cố, lập lệnh điều chuyển.
      \item \textbf{Nhà trường / Ban Quản lý:} Phân tích nhu cầu, mô phỏng kịch bản What-if, báo cáo hạ tầng.
    \end{itemize}
\end{itemize}

\subsection{Functional \& Non-functional Requirements -- Hệ thống Yêu cầu}

\textbf{Yêu cầu chức năng (Functional Requirements):}
\begin{itemize}[leftmargin=2em, itemsep=2pt]
  \item Xác định danh mục chức năng tổng thể và ranh giới nghiệp vụ của 7 phân hệ (SS1 -- SS7).
  \item Mỗi TV đặc tả chi tiết ít nhất 02 Use-case cá nhân trong phân hệ được phân công.
\end{itemize}

\textbf{Yêu cầu phi chức năng (Non-functional Requirements):}
\begin{itemize}[leftmargin=2em, itemsep=2pt]
  \item \textbf{Performance:} Thời gian phản hồi của hệ thống mô phỏng $\leq$ 3 giây cho 90\% tác vụ.
  \item \textbf{Reliability:} Hệ thống phải xử lý đúng các kịch bản biên (edge cases) và không bị crash khi nhập liệu không hợp lệ.
  \item \textbf{Availability:} Prototype Streamlit hoạt động ổn định 24/7 trong môi trường dev; không yêu cầu đăng nhập nếu chạy local.
  \item \textbf{Security:} Không lưu trữ dữ liệu cá nhân thật; dữ liệu mô phỏng được giả lập trong \texttt{st.session\-\_state} hoặc file JSON cục bộ.
\end{itemize}

\subsection{Use-case Diagram -- Sơ đồ Use-case toàn hệ thống \& Phân hệ}

\begin{itemize}[leftmargin=2em, itemsep=3pt]
  \item \textbf{System-wide Use-case Diagram (Toàn hệ thống):} Nguyễn Thành Đạt \& Phạm Tấn Đạt chủ trì xây dựng sơ đồ tổng thể (Draw.io / PlantUML) bao phủ toàn bộ 7 phân hệ, thể hiện:
    \begin{itemize}[leftmargin=1.5em, itemsep=1pt]
      \item 4 nhóm Actor chính: \textit{Sinh viên, Operator, Hệ thống Metro, Ban Quản lý / Nhà trường}.
      \item Ranh giới hệ thống (System Boundary) và các quan hệ \texttt{<<include>>}, \texttt{<<extend>>}.
      \item Lưu trữ tại: \texttt{docs/diagrams/system/system\_usecase\_diagram.[drawio|png]}.
    \end{itemize}
  \item \textbf{Individual Use-case Specification (Cá nhân):} Mỗi thành viên hoàn thành bản nháp đặc tả và vẽ sơ đồ Use-case chi tiết cho \textbf{02 Use-case cá nhân} được phân công tại Mục 3.
  \item \textbf{Cấu trúc thư mục nộp sơ đồ (\texttt{docs/diagrams/}):}
    \begin{itemize}[leftmargin=1.5em, itemsep=1pt]
      \item \texttt{docs/diagrams/system/}: Chứa sơ đồ toàn hệ thống chung của cả nhóm.
      \item \texttt{docs/diagrams/subsystems/SS<X>\_<tên\_phân\_hệ>/}: Thư mục con riêng cho từng phân hệ (SS1 đến SS7) để các cá nhân nộp bài độc lập, tránh xung đột file.
    \end{itemize}
  \item \textbf{Quy chuẩn định dạng \& Đặt tên file:}
    \begin{itemize}[leftmargin=1.5em, itemsep=1pt]
      \item Bắt buộc nộp đủ \textbf{file nguồn} (\texttt{.drawio} hoặc \texttt{.puml}) và \textbf{file ảnh xuất ra} (\texttt{.png} độ phân giải cao $\geq 300$ DPI).
      \item Cú pháp đặt tên file cá nhân: \texttt{<MSSV>\_<Mã\_Use\_Case>.[drawio|png]} \newline
      \textit{(Ví dụ: \texttt{2410709\_UC\_Dat\_Xe\_Dung\_Chung.drawio} và \texttt{.png})}.
    \end{itemize}
\end{itemize}

\subsection{Yêu cầu Nộp Nhật ký Prompt AI lên GitHub}

\begin{itemize}[leftmargin=2em, itemsep=4pt]
  \item Cuối Sprint 1, \textbf{100\% thành viên (7/7 người)} bắt buộc đẩy (push) toàn bộ nhật ký Prompt AI cá nhân lên GitHub vào đúng thư mục tên mình trong \texttt{docs/prompt-history/TVx\_HoTen/}.
  
  \item \textbf{Cấu trúc thư mục đã khởi tạo sẵn trên repository:}
    \begin{center}
    \renewcommand{\arraystretch}{1.25}
    {\small
    \begin{tabularx}{0.98\linewidth}{|>{\raggedright\arraybackslash}p{5.8cm}|X|}
      \hline
      \rowcolor{lightgray} \textbf{Thư mục trên GitHub} & \textbf{Thành viên phụ trách} \\
      \hline
      \url{docs/prompt-history/TV1_NguyenThanhDat/} & Nguyễn Thành Đạt -- SS1 \\
      \hline
      \url{docs/prompt-history/TV2_NguyenAnhTai/} & Nguyễn Anh Tài -- SS2 \\
      \hline
      \url{docs/prompt-history/TV3_PhamTanDat/} & Phạm Tấn Đạt -- SS3 \\
      \hline
      \url{docs/prompt-history/TV4_DangHoangQuyNhan/} & Đặng Hoàng Quý Nhân -- SS4 \\
      \hline
      \url{docs/prompt-history/TV5_DuongDangKhoa/} & Dương Đăng Khoa -- SS5 \\
      \hline
      \url{docs/prompt-history/TV6_HoangVanTan/} & Hoàng Văn Tấn -- SS6 \\
      \hline
      \url{docs/prompt-history/TV7_Nguyen/} & Nguyên -- SS7 \\
      \hline
    \end{tabularx}}
    \end{center}
  \item \textbf{Yêu cầu nội dung:} Chi tiết trên LMS (Link: https://lms.hcmut.edu.vn/mod/page/view.php?id=702921).
\end{itemize}

% ============================================================================
%               MỤC 5: BẢNG PHÂN CÔNG HÀNH ĐỘNG (ACTION ITEMS)
% ============================================================================
\section{Bảng Phân công Hành động (Action Items)}

\renewcommand{\arraystretch}{1.35}
\begin{tabularx}{\textwidth}{|c|X|p{4.8cm}|c|X|}
  \hline
  \rowcolor{lightgray} \textbf{STT} & \textbf{Công việc} & \textbf{Người phụ trách} & \textbf{Hạn chót} & \textbf{Tiêu chí hoàn thành} \\
  \hline
  1 &
  Khởi tạo repo Git, commit boilerplate Streamlit và tạo tài liệu kế hoạch &
  Nguyễn Thành Đạt (Leader) &
  11/09 &
  Code chạy ổn định; có \url{docs/tasks-assignment.md} trên \texttt{main} \\
  \hline
  2 &
  Khảo sát bối cảnh ĐHQG-HCM, Metro số 1; viết Mục 1.1 (Context) \& 1.2 (Stakeholders) SRS &
  Nguyễn Thành Đạt, Phạm Tấn Đạt, Đặng Hoàng Quý Nhân &
  18/09 &
  Bản thảo Context \& Stakeholders hoàn chỉnh trong SRS \\
  \hline
  3 &
  Xác định phạm vi dự án Mục 1.3 (Scope Boundary); soạn User Stories sơ bộ (3--5 stories/phân hệ) &
  Nguyễn Anh Tài, Dương Đăng Khoa \& cả nhóm &
  18/09 &
  Bản thảo Scope \& 7 bộ User Stories sơ bộ \\
  \hline
  4 &
  Chuẩn bị nội dung cho buổi Họp Tuần 2 (Review toàn bộ Phần 1) &
  Toàn thể 7 thành viên &
  19/09 &
  Tham gia họp đầy đủ lúc 20:00 ngày 20/09/2026 \\
  \hline
  5 &
  Push nhật ký Prompt AI và use-case diagram &
  Toàn thể 7 thành viên (mỗi người tự nộp) &
  19/09 &
  PR được tạo, Reviewer gán, file có timestamp \\
  \hline
\end{tabularx}

% ============================================================================
%               KHỐI CHỮ KÝ
% ============================================================================
\vspace{30pt}
\sectionrule

\vspace{10pt}

\begin{minipage}[t]{0.48\textwidth}
  \centering
  \textbf{THƯ KÝ LẬP BIÊN BẢN}\\[30pt]
  \textit{(Ký và ghi rõ họ tên)}\\[40pt]
  \dotfill\\[4pt]
  \textbf{Nguyễn Thành Đạt}
\end{minipage}
\hfill
\begin{minipage}[t]{0.48\textwidth}
  \centering
  \textbf{TRƯỞNG NHÓM PHÊ DUYỆT}\\[30pt]
  \textit{(Ký và ghi rõ họ tên)}\\[40pt]
  \dotfill\\[4pt]
  \textbf{Nguyễn Thành Đạt}
\end{minipage}

\vspace{20pt}

\begin{center}
  {\small\textcolor{darkgray}{\textit{Biên bản được lập ngày 12/09/2026, lưu trữ tại:}} \texttt{docs/meeting-minutes/}}\\[2pt]
  {\small\textcolor{darkgray}{\textit{Repository:}} \href{https://github.com/Dat-nguyen195/BTL_1_CNPM}{https://github.com/Dat-nguyen195/BTL\_1\_CNPM}}
\end{center}

\end{document}
